\section{Discussion and conclusions}
\label{sec:discussion}

The selected-distance calculation identifies a statistical coefficient under an explicit law. It provides a normalized real-row likelihood, a finite-$F$ confidence region for the joint fifteen-coordinate mean, and tensor projections that preserve nuisance uncertainty. The main limitation is not failure to propagate that uncertainty. It is the additional information needed to interpret the coefficient physically and to transport it beyond the released cohort.

The model comparison makes that limitation measurable. The vector and STF models exchange predictive ordering when the declared group covariance changes. A low-velocity object violates the perturbative interpretation of the vector response, and an earlier residual candidate fails an actual-data check. These are substantive results of model criticism. They cannot be replaced by the statement that no later fold diagnostic crosses a conservative reference threshold. The later finite-$F$ calculation changes the treatment of training uncertainty and qualifies a conditional coefficient region; it does not retrospectively erase the earlier failure or validate the parent sampling law.

Selection produces another obstruction. The exact finite-completion set retains every missing event indicator, giving a hull at least $0.56369$ wide for the 628-target fraction. Auxiliary balance identifies useful feasible weights only after the target is changed, and even then it leaves missing outcomes unresolved. This prevents the observed fraction, a balanced fraction, and a parent-population probability from being used interchangeably. Each is an estimand under different information and assumptions.

The optical-response example separates statistical from physical identification particularly clearly. Under an unrestricted additive STF bias, the identified shear image equals its declared admissible domain $\mathcal D_S$. The selected likelihood cannot distinguish its members. A bounded small-anisotropy domain therefore gives a bounded image, while unboundedness belongs only to the formal affine extension $\mathcal D_S=\STF$. A justified response-bias bound would enter the algebraic outer-support formula explicitly and could further restrict the domain-intersected image. It would still leave the other signed sectors of $x_C$ and the common reference frame to be supplied. The corresponding physical response for the separate apparent Hubble-rate carrier is presently unknown; neither this optical identified-set formula nor its formal unbounded extension is transferred between their different units.

This distinction also limits the meaning of the functional bundle. Full tensors, signs, cross blocks, and domain states carry information that scalar normalizations can remove. A reference anchor defines a comparison scale. A positive floor defines how small values are treated. Neither is inferred merely because its value makes a plot stable. In the retained four-window example, changing the effective floor moves a median ratio from $1.718693$ to $1.025396$ and eventually makes every ratio one. The observable carrier is unchanged throughout.

Fibre identification, Gaussian regression pivots, projection of a common confidence event, union coverage under a true member of a fixed family, and bounded-density event envelopes together describe the selected tensor data. Ordinary pushforward applies once a law and a defined functional are supplied. For the observed cohort, those ingredients support a conditional joint tensor region and finite-frame sensitivity sets. The physical response and the law of the missing outcomes require additional information.

Further physical inference would need source-matched response calibration, justified covariance and selection laws, and compatible frame and normalization information. A physical family comparison would additionally require family-specific predictions and null competition; a tensor amplitude or morphological match alone cannot identify a homogeneous cosmological geometry. These missing ingredients do not prevent the selected likelihood and sensitivity results from being reported. They determine their scope.

\section*{Data and computational provenance}

The observational inputs are the published Cosmicflows-4 products, the retained linewidth-target join, and identifier-matched directions from CDS/HyperLEDA. The two analysed source populations are described separately in Secs.~\ref{sec:likelihood} and \ref{sec:coupled}. The accompanying source map links each numerical table and figure to the retained results, row identities, response conventions, and scoped checks. It preserves the full coefficient and cross-window covariances rather than reconstructing them from marginal intervals. The manuscript consumes these results without new fitting or drawing. Figure preparation changes publication labels or file format while retaining the numerical content. No cross-survey likelihood multiplication or physical response imputation is used.
